% ============================================================
% 3. Method   ← 이 파일이 §3이다 (파일명은 Overleaf 템플릿 것을 그대로 둠)
% 계획: sections/03_method/README.md
%
% 구조 (2026-09-06 확정, 코퍼스 11편 조사 반영):
%   도입 = roadmap 한 문단. Notation·배경식은 두지 않는다.
%   세 소절은 "데이터셋이 아직 도착 중"이라는 한 원인의 세 결과다 (도입 주석 참조).
%   3.1 / 3.2 / 3.3  세 소절. 진단 소절은 두지 않는다.
%   진단은 각 소절 첫 문단(P1)에 분산한다 — taming3dgs §3.2 형식.
%   ※ 진단에 소절을 준 논문은 조사한 11편 중 cartgs 하나뿐이다.
%     "세 결정으로 분해"라는 우리 중심 주장은 §1 이 진다 (taming3dgs 도 그렇게 한다).
%
% \slot{move}{줄수}{할 일} = greeking. 계획 분량만큼 자리를 잡는 빈 칸.
%   초고가 들어오면 해당 \slot 을 실제 문단으로 갈아끼운다.
%   박스 안 글은 영어 (kotex 미사용). 한국어 상세는 각 slot 위 주석에.
%
% 분량 (2026-09-06 재산정). CVPR 2단 10pt, 한 쪽 ≈ 110줄.
%   ⚠ 이전 주석의 숫자(1.05p/1.0p/1.05p)는 실제 slot 합과 맞지 않았다. 아래는 실측 기준이다.
%
%   코퍼스 실측 — 같은 성격(메커니즘 명세) 소절의 총 줄수, 2단 기준:
%     CoMapGS 3.3            29 (수식2)      LM-RS 4.5              88 (수식1)
%     Taming3DGS 3.2         31 (수식1) ★    chen2026cover 4.2     101 (수식5)
%     Taming3DGS 3.4         38 (수식0)      Point-SLAM 3.1        106 (수식3)
%     LM-RS 4.3              55 (수식1)      Taming3DGS 3.3        136 (수식3)
%     SparseGS 3.2           60 (수식2)      LM-RS 4.1              67 (수식1)
%   중앙값 ≈ 63. ★ 우리 §3.1 과 문제가 가장 가까운 Taming3DGS 3.2 는 4문단 31줄로 끝낸다.
%
%   우리 배분 (수식 1개 ≈ 3줄, 그림 1개 ≈ 14줄):
%     도입  19줄 0.17p   roadmap 한 문단
%     §3.1  95줄 0.86p   문단 66 (16/12/16/10/12) + run-in 6 + 수식3 9 + Fig.3 14
%                        ※ 5문단 + run-in. taming3dgs 는 4문단이지만 그쪽은 "왜 densify 하나"를
%                          설명할 필요가 없다. 우리는 A2·A3(왜 keyframe 밖까지)가 한 문단 더 든다.
%     §3.2  101줄 0.92p  문단 84 (16/10/20/22/16) + 수식3 9
%                        ※ 최적화 문제를 display 로 올렸다. 두 항이 (a)(b) 라는 게 이 소절의 논지다.
%                        ※ P3·P4 가 각각 5가지를 해야 해서 24줄씩 준다. §3.1(0.83p)과 대등.
%     §3.3  40줄 0.36p   팀원 (채워지면 §3.1 수준으로 늘어날 것)
%   → §3 ≈ 240줄 2.18p. 팀원 §3.3 이 채워지면 ≈ 2.7p.
%   중앙값보다 조금 위에 두되 chen2026cover 4.2(101)·Point-SLAM 3.1(106) 아래로 유지한다.
%
% ============================================================
\section{Method}
\label{sec:method}

% ---------- 도입: roadmap 한 문단 (§3 도입은 이것뿐) ----------
% 묶는 논리 (2026-09-06 확정, 안 ②):
%   "데이터셋이 아직 도착 중"이라는 하나의 사실에서 세 문제가 파생된다.
%   offline 3DGS 가 딛고 선 세 가정이 그 사실 때문에 각각 깨진다:
%     (1) 고정 집합 위의 고정 예산      → 집합이 자라면 update 가 희석된다      → §3.1
%     (2) 셔플링은 유한 고정합 위에 있다 → pool 이 자라면 그 전제가 깨진다        → §3.2
%     (3) 기하는 가진 뷰로 결정된다      → free space 는 충분히 관측되기 전엔 미결정 → §3.3
%   ★ 이 틀에서 C3 는 "남은 것"이 아니라 같은 원인의 세 번째 결과다.
%   ★ §3.2 의 scope guard("pool 이 자라서 RR 보장을 못 가져온다")가 변명이 아니라
%     주제의 재확인이 된다.
%
% ⚠ 이 묶음은 carve 가 "관측 부족으로 생긴 floater"를 다룰 때 성립한다.
%   관측이 충분해도 최적화가 만든 floater 를 다루는 것이라면 이 틀은 무너지고,
%   "학습 루프의 세 손잡이(데이터/순서/목적함수)"로 되돌아가야 한다. 팀원 확인 필요.
%
% 문장 3 = 셋을 독자의 말로 먼저 세어준다 (sparsegs "consists of three key components: A, B, and C").
%   §ref 를 꺼내기 전에 "무엇을 셋으로 보는가"를 평이한 말로 심어둔다.
% 문장 4·5·6 은 같은 꼴: [깨지는 offline 가정] → [무엇을 하는가 §ref] → [무엇을 얻는가]
%   '얻는 것' 절은 lmrs 의 thereby / which results in / which eliminates.
% 코퍼스 조사: Method 도입은 5~7문장 8~13줄 한 문단이 관행.
% 원본은 plan/outline/CURRENT.md.
\slot{ROADMAP \textemdash\ the whole opening, 7 sentences}{19}{%
  (1) \emph{Premise}: online mapping fits a model to a dataset that is \textbf{still arriving} —
  new views keep entering while the model is being optimized.
  (2) \emph{Contrast}: offline 3DGS rests on three assumptions that this one fact breaks, and
  existing online systems inherit all three through a single keyframe rule.
  (3) \emph{The three points}, named in plain terms before any section reference: we therefore
  look at the three places in the mapping loop where that inheritance shows ---
  \textbf{how much} supervision may enter, \textbf{in which order} it is revisited, and
  \textbf{what constrains} the geometry it does not yet determine.
  (4) \emph{\S\ref{ssec:growth}}: \textbf{a fixed budget over a fixed set} no longer holds, so we
  let completed GPU service set how fast the set may grow, \textbf{which makes} that rate
  independent of the size of the set.
  (5) \emph{\S\ref{ssec:ercb}}: \textbf{shuffling is defined over a finite sum}, which a growing
  set is not, so we apply a bounded interval-service-shortfall correction over global
  without-replacement reshuffling, \textbf{which} prioritizes lagging coverage when one pass is
  infeasible without introducing repeated views.
  (6) \emph{\S\ref{ssec:carve}}: \textbf{geometry is determined by the views one has}, and free
  space is not yet determined --- floaters are that indeterminacy made visible --- so we
  constrain it with causal carve evidence, \textbf{which} removes them without future frames.
  (7) \emph{Figure}: Fig.~2 shows the three consequences of the one cause in a single pass of the
  mapper.}

% ---------- 3.1 (5문단, 86줄) ----------
\subsection{Compute-Paced View Growth}
\label{ssec:growth}

% 문단 구성 (2026-09-06 개정). 5문단 + run-in.
%   taming3dgs §3.2 는 4문단이지만 그쪽은 "왜 densify 하나"를 설명할 필요가 없다.
%   우리는 keyframe 밖 프레임을 쓴다는 것 자체가 설명이 필요한 선택이라 한 문단이 더 든다.
%   근거는 claims A2·A3 (둘 다 ✅).
%
%   논증: 더 쓰면 좋다(A2) → 공짜가 아니다, 최적 밀도가 예산 따라 움직인다(A3)
%         → 속도를 정해야 한다. baseline 의 답은 keyframe 이고 그건 속도를 집합 크기·
%           하드웨어에 용접한다 → 그럼 무엇이 속도를 정해야 하나 → 완료된 service(P3, 기여)
%         → 계량이 정해지면 나머지는 장부(P4, P5) → 어디에 둘지는 +0.35%(run-in, P2 로 회귀)

\begin{figure}[t]
  \centering
  \fcolorbox{slotgray}{slotbg}{\parbox[c][3.2cm][c]{0.92\linewidth}{\centering
    {\sffamily\scriptsize\bfseries\color{skelblue}FIG.\,3 \textemdash\ IN-METHOD EVIDENCE}\\[3pt]
    {\sffamily\scriptsize\color{slotgray}$x$: wall time \quad $y$: cumulative admissions\\
     law + measured + baseline (gate) \quad\textbullet\quad P01}}}
  \caption{Measured growth follows the law and is independent of the size of the training view
           set. Placed here rather than in \S\ref{sec:exp}, so the reader meets the law and its
           verification in the same place.}
  \label{fig:growthtrace}
\end{figure}

% P1 — 왜 학습 집합이 자라는가, 그리고 왜 상한을 둘 수 없는가.
%   근거: claims A2 (exp51/exp66) + neural SLAM 의 forgetting 문헌(iMAP, Point-SLAM).
%   ★ "무한정 자란다"를 여기서 세워야 §3.2 가 "그러면 어떤 순서가 이상적인가"를 물을 수 있다.
\slot{P1 \textemdash\ WHY THE SET GROWS, AND WHY IT CANNOT BE CAPPED}{16}{%
  Keyframes are selected for tracking, not for photometric supervision. Under the same objective a
  view set that includes frames between keyframes reaches a lower held-out loss, and those frames
  have \emph{already arrived}: nothing is spent to acquire them \emph{(claims A2, exp51/exp66)}.
  \textbf{Nor can the set be capped by a sliding window.} Evicting a view stops supervision for
  the region it covers, and the parameters are shared, so continued optimization elsewhere
  degrades it --- this is why keyframe replay exists in neural SLAM at all [iMAP, Point-SLAM].
  \textbf{The training set therefore grows monotonically with the trajectory}, which is the
  premise \S\ref{ssec:ercb} has to schedule against.}

% P2 — 왜 keyframe 규칙이 틀린 답인가. 근거: claims A3 (exp66 viewset density).
%   ⚠ A2 의 경계를 여기서 밝힌다: dense 를 쓴다는 것 자체는 전제이지 우리 주장이 아니다.
\slot{P2 \textemdash\ WHY THE KEYFRAME RULE IS THE WRONG ANSWER}{12}{%
  Using them is not free, however: every admitted view competes for the same fixed number of
  optimizer updates, so the best view-set density moves with the compute budget
  \emph{(claims A3, exp66)}. A rate must therefore be chosen. The baseline chooses it by keyframe
  --- ``For each new keyframe, we run 10 mapping iterations'' --- which welds that rate to the
  size of the set and to hardware speed. \textbf{Say here that using dense supervision is the
  premise, not the claim} \emph{(boundary of A2)}; what we claim is the rule that sets the rate.}

% P3 — 이 절의 기여. 세 후보를 비교하고 계보를 밝힌다.
%   ⚠ 2026-09-08: 단일 GPU·pace 스윕 없음으로 정하면서 "하드웨어 속도와 무관" 주장은
%     검증할 실험이 사라졌다(claims B2, Fig.4, P02 폐기). 집합 크기 독립성만 말한다 —
%     그건 한 run 안에서 pool 이 커지는 동안 측정된다(B1, exp73 526 poll).
%     근거: notes/decisions/2026-09-08_single-gpu-realtime.md
\slot{P3 \textemdash\ WHAT TO METER (the contribution)}{16}{%
  Three candidates. \emph{Wall-clock time} advances whether or not the GPU finished anything, so
  under a slow or stalled optimizer admission outruns it and the new views dilute the updates ---
  open loop. \emph{Iteration count} is closer, but the cost of an iteration grows with the set,
  so metering iterations keeps the coupling we are trying to remove. \emph{Completed GPU service}
  is the only closed loop: it advances only when the optimizer does. Lineage: this is a token
  bucket [RFC 2212] whose tokens are minted by completed service rather than by elapsed time,
  which makes it credit-based rather than rate-based. A consequence: during a stall the meter
  does not advance, so no credit accrues and no burst can build --- a time-metered bucket needs a
  depth cap for that, ours cannot fill in the first place.
  \pend{claim set-size independence only. "hardware speed" dropped 2026-09-08: single GPU, no
  pace sweep, so B2 has no experiment --- notes/decisions/2026-09-08\_single-gpu-realtime.md}}

% P4 — 불변식이 먼저, 상한이 그로부터 따라온다.
\slot{P4 \textemdash\ INVARIANT, AND THE CAP THAT FOLLOWS}{10}{%
  One invariant: credit spent may never exceed credit earned, Eq.~(\ref{eq:invariant}).
  Solving it for the number of admissions gives the cap, Eq.~(\ref{eq:astar}) --- the floor
  because admissions are integral, the $\min$ because views that have not arrived cannot be
  admitted. Two free parameters only: $\kappa$, the cost of one admission, and $\gamma$, the
  conversion from service to credit. \pend{$\kappa$ was fixed on one scene and reused; say so}
  \pend{decide: is (\ref{eq:astar}) stated on cumulative $S(t)$, or on credit remaining after the
  discards of P5? The two differ once a discard fires, and a reviewer will check}}

\begin{equation}\label{eq:invariant}
  \kappa\big(|A_t| - A_0\big) \;\le\; B_0 + \gamma S(t)
\end{equation}
\begin{equation}\label{eq:astar}
  A^{*}(t) \;=\; \min\!\Big(M(t),\; A_0 + \Big\lfloor \tfrac{B_0+\gamma S(t)}{\kappa} \Big\rfloor \Big)
\end{equation}

% P5 — 현실이 어긋나는 두 지점 + 정책 선언 + forward pointer.
%   ⚠ claims B5: no-prepurchase 는 구현된 정책이다. carry 보다 낫다는 주장은 금지.
\slot{P5 \textemdash\ REALIZED QUOTA, DISCARD, POLICY}{12}{%
  Reality diverges twice. Only views that have arrived and are not yet admitted are candidates,
  giving the realized quota Eq.~(\ref{eq:quota}). And when nothing is pending, credit is
  \emph{discarded} rather than banked --- a second guard against bursts, and what keeps admission
  causal, since credit never crosses an interval boundary to be spent on views that arrive later.
  The implemented policy is therefore \emph{no-prepurchase}.
  \pend{report it as the policy we implement, NOT as one shown better than carry --- B5}
  \S\ref{sec:exp} shows the realized count matches Eq.~(\ref{eq:astar}) with zero integer error,
  and how quality degrades as $\kappa$ falls (Fig.~\ref{fig:growthtrace}).}

\begin{equation}\label{eq:quota}
  Q_t \;=\; \min\!\big(|C_t|,\; [\,A^{*}(t) - |A_{t^-}|\,]_+\big)
\end{equation}

% run-in — 기여 아님. 약하다는 사실 자체가 P2 의 논거를 되받는다.
\noindent\textbf{Slot placement.}\;
\slot{RUN-IN \textemdash\ NOT A CONTRIBUTION}{6}{%
  The quota says how many; temporal maximin and interval water filling say where along the
  trajectory. The measured gain is $+0.35\%$ --- the bottleneck is \emph{how many} views can be
  admitted, not \emph{which} ones, which is why this subsection is about the rate.}

% ---------- 3.2 ----------
% 2026-09-15 authority correction: the old per-view raw-count Gibbs prototype
% exp(-beta n_i) is not the production/paper-Full method.  The source of truth
% is Stage 3d: interval completed-service shortfall, a bounded Gibbs bonus,
% transactional accounting, and a global without-replacement residue.
\subsection{Service-Shortfall Entropy-Regularized Count Balancing}
\label{ssec:ercb}

\slot{P1 \textemdash\ THE REGIME AND THE SCHEDULING FAILURE}{15}{%
  Random reshuffling is already sufficient when every admitted view can receive at least one
  service before the map is read.  The failure appears in the opposite regime.  Let $U$ be the
  number of unique views eligible before a dense-service opportunity and $S$ the number of
  completed dense-view Adam services.  If $S<U$, one global pass is infeasible: a growing pool can
  leave whole causal intervals untouched even though the order inside the serviced subset is
  random.  We therefore correct \emph{interval service shortfall}, rather than trying to equalize
  every view's lifetime count.  This is a conditional claim about structurally scarce service, not
  a claim that the correction must outperform reshuffling when $S\ge U$.}

\slot{P2 \textemdash\ COMPLETED-SERVICE SHORTFALL}{18}{%
  Partition the currently admitted causal set $A_t$ by the keyframe bracket known when each view
  was admitted.  For interval $G_j$, let $m_j=|G_j\cap A_t|$ and let $c_j$ count only selections
  whose Adam update has completed in the current map generation.  Its per-member service rate is
  $r_j=c_j/m_j$, and the frame-weighted global mean is
  $\mu=\sum_j c_j/\sum_j m_j$.  Equation~(\ref{eq:ercb-shortfall}) measures only shortage below a
  relative floor $\rho\mu$; when the target is zero we set every deficit to zero.  Hence
  $d_j\in[0,1]$, the statistic is invariant to a common scaling of all completed counts, and an
  interval at or above the floor receives no correction.}

\begin{equation}\label{eq:ercb-shortfall}
 d_j=\begin{cases}
  \left[1-\dfrac{r_j}{\rho\mu}\right]_+, & \rho\mu>0,\\[2pt]
  0, & \rho\mu=0.
 \end{cases}
\end{equation}

\slot{P3 \textemdash\ BOUNDED ENTROPY-REGULARIZED MEASURE}{17}{%
  Let $E_t\subseteq A_t$ be the unserved residue of the current global reshuffling epoch and
  $a_j=|G_j\cap E_t|$.  At an outer-block boundary, interval $j$ receives the frame-uniform base
  mass $a_j$ multiplied by the bounded Gibbs bonus in Eq.~(\ref{eq:ercb-weight}).  The factor
  $a_j$ preserves uniform-per-frame reshuffling when $\gamma=0$; the exponential term changes only
  the interval order and is bounded by $1\le\exp(\gamma d_j)\le\exp(\gamma)$.  We use
  $\rho=0.75$ and $\exp(\gamma)=1.5$, fixed before the cross-sequence confirmation.}

\begin{equation}\label{eq:ercb-weight}
 w_j=a_j\exp(\gamma d_j),\qquad
 p_t(j)=\frac{w_j}{\sum_{\ell:a_\ell>0}w_\ell}.
\end{equation}

\slot{P4 \textemdash\ BLOCK DRAW AND GLOBAL RESIDUE}{19}{%
  We draw up to $K=8$ distinct intervals from $p_t$ without replacement using a
  Plackett--Luce (equivalently, Gumbel-Top-$k$) block.  A persistent shuffled inner queue chooses
  one member of $G_j\cap E_t$ for each selected interval.  Crucially, without replacement is
  global rather than merely interval-local: a committed view leaves $E_t$, and no view may repeat
  until the current residue is exhausted.  Already-arrived candidates may join the residue while
  an epoch grows, but cannot enter an outer block that has already been materialized.  When
  $E_t$ becomes empty, the next draw starts a new epoch from the then-current $A_t$.  This retains
  reshuffling's coverage invariant while allowing a bounded reordering toward lagging intervals.}

\slot{P5 \textemdash\ TRANSACTION AND CAUSALITY CONTRACT}{17}{%
  Selection is a proposal, not service.  The scheduler increments $c_j$ and removes a view from
  $E_t$ only after the owning Adam step commits; a deadline rejection cancels the proposal and
  restores the same concrete block without advancing the service clock.  Interval identity is
  frozen from already-observed admission endpoints, so later pose-graph refreshes cannot rewrite
  scheduler history.  A map reset clears the service counts and residue.  Thus C1 controls which
  views enter, whereas ERCB changes only the order in which the shared eligible set is serviced.
  We report rich sequences as controls and test the ERCB advantage only on workload-ledger-defined
  $S<U$ sequences; no image-quality metric is used to choose that label.}

% ---------- 3.3 (팀원, 110줄) ----------
\subsection{Causal Free-Space Carving}
\label{ssec:carve}

\slot{\color{red}P1--P4 \textemdash\ TEAMMATE}{40}{%
  Match the shape of the two subsections above:
  Diagnosis $\to$ Mechanism $\to$ assumption ledger $\to$ forward pointer to
  \S\ref{sec:exp} (\texttt{Carving.}).
  \textbf{Use the same name announced in P3 of the section opening.}
  The subsection title is the teammate's to fix.}
